\begin{textsample}{NEG Dim 2 – 2010\_c5 – Score -2.00 – 2010\_c5/enade\_2010\_tec\_agro\_87acf19614.txt}  \label{ex:f2_neg_103}
CAPÍTULO 5 DISTRIBUIÇÃO DOS CONCEITOS 5.1 PANORAMA NACIONAL DA DISTRIBUIÇÃO DOS CONCEITOS A Tabela 5.1 apresenta a quantidade e distribuição de cursos de Tecnologia em Agronegócio participantes do ENADE/2010 por faixa de conceito e Grande Região.
Observando -se os dados da Tabela 5.1, nota -se que, dos 57 cursos participantes, cerca de 25,0 \% ( 14 cursos ) classificaram -se com conceito 3, o valor modal, o que não representa um padrão nacional.
O segundo conceito mais atribuído foi o 4, presente em quatro das cinco regiões.
Dos 57 cursos, 24 ( 42,1 \% ) ficaram sem conceito ( SC ).
A região Sul participou com 21 cursos, 36,8 \% do total nacional.
Nesta região, o conceito 3 foi atribuído a seis cursos dos 21 participantes, equivalente a 28,6 \% do total regional.
Um pouco mais da metade dos cursos ( 52,3 \% correspondendo a 11 cursos ) não receberam conceito, outros dois cursos ( 9,5 \% ) receberam conceito 5 e os dois restantes receberam conceitos 2 e 4.
A região Sudeste participou com 17 cursos, 29,8 \% do total nacional.
Destes, cinco cursos, 29,4 \% em termos regionais, obtiveram conceito 3, o modal.
O conceito 1 foi atribuído a 5,9 \% dos cursos ( um curso ), o conceito 2 a três cursos ( 17,6 \% ), o conceito 4 foi atribuído a dois cursos ( 11,8 \% ) e o conceito 5 a um curso.
Foram cinco os cursos sem conceito ( 29,4 \% ).
A região Nordeste, por sua vez, participou com apenas dois cursos, mas estes não receberam conceito.
Dos 12 cursos participantes da região Centro-Oeste, quatro ( 33,3 \% ) obtiveram conceito 2, o conceito modal.
O conceito 3 foi atribuído a dois cursos e outros quatro ficaram sem conceito.
Os conceitos 4 e 5 foram atribuídos a um curso cada ( 8,3 \% ).
Na região Norte a participação de cinco cursos foi observada, mas apenas três receberam conceito, sendo dois com conceito 4 ( 40,0 \% ) e um com conceito 3 ( 20,0 \% ).
Os outros dois cursos ficaram sem conceito ( 40,0 \% do total regional ). 5.2 CONCEITOS POR CATEGORIA ADMINISTRATIVA E POR GRANDE REGIÃO A Tabela 5.2 apresenta a distribuição dos cursos participantes do ENADE/2010 de Tecnologia em Agronegócio, por Categoria Administrativa, de acordo com os conceitos por eles alcançados, segundo as Grandes Regiões brasileiras.
Dos 57 cursos participantes, 12 ( 21,1 \% ) eram ministrados em instituições públicas e 45 ( 78,9 \% ), em privadas.
De acordo com as informações da Tabela 5.2, as instituições privadas foram as que apresentaram, em termos nacionais, a maior proporção de cursos com conceito 3, o conceito modal, equivalente a 26,7 \% ( 12 cursos ) do total dessa categoria.
Na rede pública, o conceito modal ficou entre 3 e 5 ( com dois cursos cada ).
Entre os demais cursos participantes, um recebeu conceito 2 e um recebeu conceito 4.
Nesta categoria, seis cursos ficaram sem conceito.
Na análise por região, observa -se que, na região Sul, as instituições privadas participaram com 16 cursos, dos quais sete ficaram sem conceito.
O conceito modal para as instituições privadas na região Sul foi 3 ( com cinco cursos ) e os restantes foram avaliados com conceito 4 ( um curso ), conceito 2 ( um curso ) e conceito 5 ( dois cursos ).
As instituições públicas participaram com cinco cursos, dos quais quatro ficaram sem conceito.
O conceito modal para esta Categoria Administrativa na região Sul foi 3 ( com um curso ).
Na região Sudeste, a rede privada concentrou 15 dos 17 cursos participantes, 88,2 \% do total da região.
Entre os que receberam conceito nesta Categoria Administrativa na região, \textbf{predominaram} os cursos com conceito 3 ( quatro cursos ) – o conceito modal.
Os demais foram classificados com conceito 1 ( um curso ), conceito 2 ( três cursos ) e conceito 4 ( dois cursos ), ou ficaram sem conceito ( cinco cursos ).
As instituições públicas, nesta região, participaram com dois cursos, sendo este avaliados com conceitos 3 e 5.
Na região Centro-Oeste, a proporção de cursos da rede privada, 83,3 \%, foi a segunda mais elevada em relação às demais regiões brasileiras, correspondendo a dez dos 12 cursos participantes.
Nesta categoria na região Centro-Oeste, o conceito modal foi 2 ( três cursos ).
Os demais cursos foram avaliados com conceito 3 ( dois cursos ), conceito 4 ( um curso ) e quatro cursos ficaram sem conceito.
Nessa região as instituições públicas participaram com dois cursos, sendo estes avaliados nos conceitos 2 e 5.
As instituições privadas concentraram três dos cinco cursos participantes da região Norte, 60,0 \% do total regional.
Desses, um ficou sem conceito e os outros dois receberam os conceitos 3 e 4.
Dos dois cursos em instituições públicas, um ficou sem conceito e o ourro recebeu conceito 4.
Na região Nordeste, foram apenas dois cursos participantes, sendo um em instituição pública e o outro em instituição privada.
Os dois cursos participantes receberam, tanto em IES pública como em privada, o conceito 5. 5.3 CONCEITOS POR ORGANIZAÇÃO ACADÊMICA E POR GRANDE REGIÃO Na Tabela 5.3 encontra -se a distribuição dos conceitos atribuídos aos cursos participantes do ENADE/2010 na Área de Tecnologia em Agronegócio, por Organização Acadêmica, segundo as Grandes Regiões brasileiras.
Dos 57 cursos participantes, 19 ( 33,3 \% ) eram oferecidos em Universidades, oito ( 14,0 \% ) em Centros Universitário e os demais 30 ( 52,7 \% ) em Faculdades.
De acordo com os dados apresentados, a maioria dos cursos avaliados com conceito 3 eram vinculados a Faculdades ( dez dos 14 cursos com este conceito ), o conceito modal.
Além dos dez cursos que ficaram sem conceito nesta Categoria Administrativa, seis receberam conceito 2 e os outros quatro receberam conceitos 4 e 5 ( com dois cursos cada ).
Entre os cursos em Centros Universitários, o conceito modal ficou entre os valores 3 e 4 ( com dois cursos cada ), e dois ficaram sem conceito.
Os outros dois cursos neste tipo de Organização Acadêmica receberam os conceitos 1 e 2.
Nas Universidades, dos 19 cursos, 12 ficaram sem conceito e um recebeu conceito 2.
Os outros seis cursos receberam os conceitos 3, 4 e 5 ( com dois cursos cada ).
Considerando -se as regiões brasileiras, verifica -se que, nas regiões Nordeste e Sul, 50,0 \% dos cursos estavam em Universidades.
A região Norte, por sua vez, concentrou 60,0 \% dos cursos em Faculdades.
Mas foi na região Sudeste que se deu a maior concentração de cursos em Centro Universitários ( 29,4 \% ).
Ao analisar separadamente cada região foi possível observar que, na região Nordeste, as Universidades participaram com um curso dos 2 cursos na Área de Tecnologia em Agronegócio da região e este ficou sem conceito.
O outro curso participante era desenvolvido em Faculdade, mas também ficou sem conceito.
Na região Sudeste as Faculdades concentraram nove dos 17 cursos.
Entre os cursos em Faculdades na região Sudeste, o conceito modal ficou entre 2 e 3 ( com dois cursos cada ), e três ficaram sem conceito.
Os demais cursos receberam conceito 4 ( um curso ) e conceito 5 ( um curso ).
Os Centros Universitários participaram com cinco cursos na região, dos quais dois obtiveram o conceito modal, 3, um ficou com conceito 2, um ficou com conceito 1 e o último ficou sem conceito.
As Universidades foram representadas por três cursos na região Sudeste, que se \textbf{distribuíram} nos conceitos 3 ( um curso ), 4 ( um curso ) e um cursos ficou sem conceito.
Dos 21 cursos da região Sul, dez eram de Universidades, sendo que oito ficaram sem conceito.
Os outros dois cursos receberam, respectivamente, os conceitos 3 e 5.
Os Centros Universitários tiveram participação de apenas um curso e este recebeu conceito 4.
As Faculdades, por sua vez, participaram com dez cursos na região, sendo que três não receberam conceito.
O conceio modal nesta categoria de Organização Acadêmica foi 3 ( com cinco cursos ) e os outros dois cursos receberam os conceitos 2 e 5.
Na região Centro-Oeste, três dos 12 cursos eram de Universidades.
Nesse tipo de organização, um curso recebeu conceito 5, um curso recebeu conceito 3 e o outro ficou sem conceito.
Nos Centros Universitários a participação foi de dois cursos na região Centro-Oeste, sendo que um ficou sem conceito e o outro recebeu o conceito modal, 4.
Nas Faculdades a participação de cursos na região foi maior, sete.
O conceito modal, 2, foi atribuído a três cursos, o conceito modal e outros dois ficaram sem conceito.
Os dois cursos restantes receberam o conceito 3.
A região Norte não apresentou cursos em Centros Universitários.
A concentração se deu em cursos de Faculdade três dos cinco oferecidos.
Destes três cursos, um não recebeu conceito e os outros dois receberam, respectivamente, os conceitos 3 e 4.
As Universidades, por sua vez, participaram na região com dois cursos, sendo que um ficou sem conceito e o outro recebeu conceito 4.

% matched lemmas: distribuir, predominar
\end{textsample}
